% !TEX root = ../Main.tex
\chapter{切线}

\section{切线与曲线的交点个数}

初中学圆的切线时，"切线与圆只有一个公共点"。但对一般曲线，这句话\textbf{不一定成立}——切线与曲线的交点个数不一定是 $1$ 个。下图里那条直线在某点与曲线相切，可在别处又穿过曲线，一共有 $2$ 个交点：

\begin{center}
\begin{tikzpicture}[>=Stealth,scale=1.0]
  \draw[thick,blue!70!black,domain=-1.3:2,smooth,samples=80]
      plot (\x,{0.6*\x*\x*\x-1.0*\x});
  \draw[thick] (-1.6,-0.5) -- (2.2,-0.5);
  \fill (-0.91,-0.5) circle (1.3pt);
  \fill (1.49,-0.5) circle (1.3pt);
  \node at (1.7,-0.9) {\footnotesize 2 个交点};
\end{tikzpicture}
\end{center}

可见"只有一个公共点"不能用来定义一般曲线的切线，得换一个刻画。

\section{从变化率的角度再看切线}

\state{切线（变化率刻画）}{
    切线是过曲线上一点、\textbf{且能刻画该曲线在该点变化率}的直线。
}

也就是说，切线的斜率就是曲线在切点处的变化率（这正是后面"导数 = 切线斜率"的来源）。为什么"能刻画变化率"就该是切线？可以用\textbf{反证法}说明。

设 $\ell$ 过曲线上的点 $A$，若 $\ell$ 的斜率 $k_{AB'}$ 与曲线在 $A$ 处的真实变化率 $k_{AB}$ 不相等：
\[
k_{AB'}\ne k_{AB},
\]
那么 $\ell$ 与曲线在 $A$ 附近就会"岔开"，必然与曲线再相交于另一点 $B'$（如下图，$\ell=AB'$ 是穿过曲线的割线，而真正贴着曲线的是 $AB$ 方向）。这样的 $\ell$ 只是割线，从而\textbf{不能}成为曲线在 $A$ 处的切线。因此切线的斜率必须等于该点的变化率。

\begin{center}
\begin{tikzpicture}[>=Stealth,scale=1.1]
  \draw[thick,blue!70!black,domain=-1.1:1.1,smooth,samples=80]
      plot (\x,{\x*\x});
  \coordinate (A) at (0.55,0.3025);
  % tangent direction at A: slope 2*0.55=1.1
  \draw[red!70!black,thick] ($(A)+(-0.75,{-0.75*1.1})$) -- ($(A)+(0.55,{0.55*1.1})$)
      node[right]{\footnotesize $AB$（切）};
  % secant AB'
  \draw[thick] ($(A)+(-0.95,{-0.95*0.45})$) -- ($(A)+(0.5,{0.5*0.45})$)
      node[right]{\footnotesize $AB'$（割）};
  \fill (A) circle (1.3pt) node[below right]{$A$};
  \fill (-0.62,0.3844) circle (1.1pt) node[above left]{$B'$};
\end{tikzpicture}
\end{center}
